\chapter{Deployment}
\label{chap:ch4}

\par When it comes to deploying an application, there are many paths one could take. One might be ordering a physical server at home and running your application on there, another one could be running it locally on your machine where only people in your network can access it, and another one could be deploying the application in the cloud. The last one was chosen for Doclane.

\section{Cloud computing}

\par Cloud computing was preferred since it is the more modern approach, and reduces overhead with managing physical servers. Without the burden and expenses of buying, setting up and maintaining servers, more effort could be redirected to the development of the app and adding features that bring more value to it. It is also that renting those servers means you can just get rid of them whenever you're not using them anymore, and you only pay for how much you use them.

\section{Amazon Web Services}

\par The deployment of the app was done on AWS. The architecture leverages containers orchestrated by Kubernetes, specifically the managed EKS (Elastic Kubernetes Service) offering from AWS.

\par AWS is the leading cloud provider, covering roughly 31\% of the global market as of 2026. Part of the reason it is so widely used is because it was the first to market, so it had the most amount of time out of all competitors to refine their services and offerings, which is half of the reason it was chosen. The other half was because the author of this thesis project is a certified Solutions Architect Professional on AWS and had a bias to pick this provider.

\subsection{Containers and Kubernetes}

\par A container is a lightweight, standalone package that includes everything needed to run a piece of software: the code, runtime, libraries and dependencies. Unlike virtual machines, containers share the host operating system kernel, making them fast to start and efficient with resources.

\par Docker is the tool used to build these containers. The backend and frontend of Doclane each have a Dockerfile that defines how the container image is built. The backend produces a small static Go binary running on a minimal base image, while the frontend uses Next.js standalone output on a Node.js base image.

\par Kubernetes is an open-source container orchestration platform. It takes your containers and runs them across a cluster of machines, handling scaling, networking, health checks and restarts automatically. Instead of manually starting containers on servers, you declare what you want (for example: "run 2 copies of the backend") and Kubernetes makes it happen.

\par EKS (Elastic Kubernetes Service) is AWS' managed Kubernetes offering. AWS handles the control plane (the brain of the cluster), while you manage the worker nodes (the machines that actually run your containers). This removes the operational burden of maintaining Kubernetes itself while giving full access to its features.

\section{The architecture}

\par The architecture consists of an EKS cluster running in a VPC with two availability zones. Traffic enters through an ALB (Application Load Balancer) that routes requests to the appropriate service: the backend or the frontend. The backend connects to an RDS PostgreSQL database and an S3 bucket for document storage. Authentication is handled by Cognito.

\subsection{Networking}

\par The networking part of Doclane consists of a VPC (Virtual Private Cloud) in the eu-west-1 Region on AWS. This VPC consists of two private subnets and two public subnets. The private subnets host the EKS worker nodes and the RDS database, so that no direct public internet access can be established to them.

\par The public subnets host the NAT Gateway and the ALB (Application Load Balancer). The NAT Gateway allows resources in private subnets to reach the internet (for example to pull container images) without being directly accessible from the internet.

\par Both the private and public subnets span two AZs (Availability Zones), such that high availability is achieved. In case one data center from AWS goes down where our application runs, then it will still continue to run as if nothing has happened.

\par An internet gateway was provisioned and added to the route table of the public subnet, such that internet access is possible to and from it. Private subnets are associated with a route table that routes internet-bound traffic through the NAT Gateway. So the thing that makes a subnet public or private is whether or not the associated route table has a direct route to the internet gateway or goes through a NAT Gateway.

\par The subnets are tagged with Kubernetes-specific labels so that the AWS Load Balancer Controller can automatically discover which subnets to place ALBs in.

\subsubsection{Security}

\par When it comes to network security, the architecture leverages Security Groups for securing traffic at the pod level and at the RDS level.

\par A Security Group in AWS is a stateful firewall that operates at a single resource level. It consists of inbound / outbound rules, inbound rules meaning what kind of traffic is allowed (on which port and from which source), while outbound rules control what traffic is allowed to go out of our resource.

\par The RDS database is protected by a Security Group that only allows inbound traffic on port 5432 from the EKS cluster Security Group. This means only pods running in the cluster can reach the database. Any other traffic is dropped.

\par The entire communication between the backend pods and RDS happens over encrypted connections (SSL/TLS is enforced on the database). Traffic between pods within the cluster uses the internal Kubernetes network and never leaves the VPC.

\subsection{Database layer}

\par RDS (Relational Database Service) was the preferred choice for the database. As a relational database was needed, the natural choice was AWS' managed service for this.

\par RDS is built in such a way that AWS handles all the operating system patching of the server that is running your database, and you also benefit from automated backups out of the box.

\par The engine is PostgreSQL, chosen for being popular, fast and well supported by RDS. The database schema is automatically applied on every fresh deployment by a Kubernetes Job that runs the initialization SQL script.

\subsection{Backend layer}

\par The backend runs as a Deployment in Kubernetes with 2 replicas by default. Each replica is a container running the Go binary that serves the chi HTTP router on port 8080. 

\par A HorizontalPodAutoscaler is configured to scale the backend from 2 to 6 replicas based on CPU utilization. If average CPU goes above 70\%, Kubernetes automatically adds more pods to handle the load.

\par Health checks are configured using Kubernetes liveness and readiness probes. The liveness probe ensures crashed containers get restarted. The readiness probe ensures traffic is only sent to pods that are ready to serve requests.

\subsubsection{IAM Roles for Service Accounts (IRSA)}

\par The backend pods need access to AWS services like S3, Textract, Bedrock, Polly and Cognito. Instead of using static credentials, the pods use IRSA (IAM Roles for Service Accounts).

\par IRSA works by linking a Kubernetes Service Account to an IAM Role. When a pod uses that Service Account, the AWS SDK automatically receives short-lived credentials scoped to exactly the permissions defined in the IAM Role. No static keys exist anywhere, not in the image, not in environment variables, not on disk.

\par The IAM policy attached to the backend pod role follows the principle of least privilege. It only allows the specific actions needed: uploading and downloading from the documents S3 bucket, calling Textract for text extraction, calling Bedrock for document interpretation, calling Polly for speech synthesis, and looking up users in Cognito.

\subsection{Frontend layer}

\par The frontend runs as a Deployment in Kubernetes with 2 replicas. Each replica runs the Next.js standalone server on port 3000. The standalone output mode produces a minimal server bundle that does not need the full node\_modules directory, keeping the container image small.

\par A HorizontalPodAutoscaler scales the frontend from 2 to 4 replicas based on CPU utilization.

\par Server-side rendering (SSR) calls from the frontend to the backend happen over the internal Kubernetes network using the service DNS name. This means SSR requests never leave the cluster and have minimal latency.

\subsection{Ingress and Load Balancing}

\par Traffic enters the cluster through an ALB (Application Load Balancer) that is automatically provisioned by the AWS Load Balancer Controller. This controller runs inside the cluster and watches for Ingress objects. When it sees one, it creates and configures a real ALB in AWS to match the routing rules defined in the Ingress.

\par The Ingress defines two routing rules: requests to /api are sent to the backend service, and all other requests are sent to the frontend service. The ALB terminates TLS using an ACM certificate and redirects HTTP traffic to HTTPS.

\par The ALB uses IP-mode target groups, meaning it routes traffic directly to pod IPs rather than going through node ports. This is more efficient and allows for faster health checks.

\subsection{DNS and HTTPS}

\par The application is served at thesis.robert-beres.com. An ACM (AWS Certificate Manager) certificate was provisioned for this domain and validated automatically via DNS. The certificate is attached to the ALB so all traffic is encrypted.

\par A Route53 alias record points the domain to the ALB. When the infrastructure is torn down and recreated, the ALB gets a new DNS name, but Terraform automatically updates the Route53 record to point to the new one.

\subsection{Document storage}

\par S3 (Simple Storage Service) is used for storing uploaded documents. The bucket has versioning enabled so previous versions of files are preserved, and server-side encryption is applied to all objects.

\par The backend generates presigned URLs for downloads. This means clients download directly from S3 using a temporary URL that expires after a set time, without any permanent credentials ever existing. This makes the process so much more secure since if somehow credentials get leaked the attack surface is very small due to the scoped permissions, and the attack surface is also ephemeral.

\subsection{Authentication}

\par Cognito is used as the identity provider. It manages user registration, login, password policies and token issuance. The frontend uses the Amplify library to interact with Cognito directly from the browser.

\par After a user logs in, Cognito issues a JWT (JSON Web Token). This token is stored in an HTTP-only cookie and sent with every request. The backend validates the token by checking its signature against Cognito's public keys (JWKS).

\subsection{Logging and monitoring}

\par Container logs from all pods are collected by Kubernetes and can be viewed using kubectl. The structured JSON logging from the backend makes it easy to filter and search for specific events or errors.

\section{Machine Learning}

\par This section goes into detail on the AWS services leveraged to build the machine learning assisted features aforementioned in the implementation chapter.

\subsection{Textract}

\par Textract is a fully managed serverless service offered by AWS that extracts text from documents. API calls are made to Textract from Doclane in order to extract text from documents.

\subsection{Polly}

\par Polly is a fully managed serverless service offered by AWS that generates speech from text. API calls are made to Polly to generate speech from selected documents.

\subsection{Bedrock}

\par Bedrock is a fully managed serverless service offered by AWS that gives you access to different LLMs that you can prompt. In our case we use the Nova Lite model from Amazon to interpret whether a document provided by an user is actually what it is required by the request.

\section{Infrastructure as Code - Terraform}

\par Whenever deploying applications in the Cloud, you can manually click in the AWS console and provision everything you need. It's straight forward, intuitive and it provides a friendly experience. Although this is convenient, it is a manual process. And manual processes are prone to errors. It becomes easy to forget what change in configuration you made last week, more simply put, it's a mess.

\par That's where the concept of IaC (Infrastructure as Code) comes into play. This involves specifying your infrastructure in a declarative manner, and having everything defined as code, it becomes much easier to maintain and have control over.

\par Terraform by HashiCorp is the chosen tool for the project. Unlike CloudFormation which is AWS-specific, Terraform is cloud-agnostic and has a large ecosystem of providers. It was chosen for its ergonomics: features like force\_destroy on S3 buckets, native Helm and Kubernetes providers, cross-region resource management with provider aliases, and the ability to manage the entire stack (AWS resources, Kubernetes manifests, Helm charts) from a single tool.

\subsection{State management}

\par Terraform keeps track of what it manages in a state file. This file maps the resources defined in code to the real resources in AWS. It is stored remotely in an S3 bucket with versioning enabled, so it can be recovered if corrupted. A DynamoDB table is used as a lock to prevent two people (or two CI runs) from modifying infrastructure at the same time.

\subsection{Stack layout}

\par The infrastructure is split into multiple Terraform stacks, each with its own state file:

\par \textbf{bootstrap} - Creates the S3 bucket and DynamoDB table for state management. Run once, never destroyed.

\par \textbf{platform-data} - Contains long-lived resources that are cheap to keep running: ECR container registries, Cognito user pool, ACM certificate, GitHub OIDC role. These persist across compute teardowns.

\par \textbf{platform-compute} - Contains the expensive resources: VPC, NAT Gateway, EKS cluster, RDS database, S3 documents bucket, IRSA roles. These are created when working and destroyed when not needed.

\par \textbf{workload} - Contains everything inside the Kubernetes cluster: Helm charts, deployments, services, ingress, database seeding jobs. Destroyed together with platform-compute.

\par This split allows tearing down the expensive compute resources while not using the app, since in development you do not need the app running 24 hours a day, while keeping the cheap persistent resources alive. A full environment can be recreated from scratch with a single command.

\subsection{Ephemeral environments}

\par The entire compute layer is designed to be ephemeral. Running \texttt{make up} provisions the full environment from scratch in about 15 minutes: VPC, NAT, EKS cluster, RDS database, container deployments, database schema seeding, admin user creation, ALB, HTTPS, DNS record. Running \texttt{make down} destroys everything in about 10 minutes, bringing the cost to zero.

\par This is possible because:
\begin{itemize}
    \item Container images persist in ECR between sessions
    \item The Cognito user pool persists so users do not need to re-register
    \item The ACM certificate persists so DNS validation does not need to happen again
    \item The database schema is applied automatically by a Kubernetes Job on every fresh deployment
    \item The admin user is bootstrapped automatically by another Kubernetes Job
\end{itemize}

\section{Continuous Integration / Continuous Deployment}

\par Modern practices of deployment have been adopted, such that the deployment pipeline of the application is fully automated. Once changes are pushed to the main branch, container images are automatically built and pushed to ECR.

\par CI/CD has been done with the help of GitHub Actions. GitHub Actions are convenient to use when your repository is already on GitHub, hence why they were chosen.

\par The pipeline builds Docker images for the backend and frontend, tags them with the git commit SHA for traceability, and pushes them to ECR. The backend and frontend are built independently, only when their respective code changes.

\par Updates happen leveraging the least-privilege principle, without any permanent credentials ever existing or being stored anywhere. The GitHub Action worker generates an ID token using GitHub's OIDC (OpenID Connect). GitHub signs this token with their private key, and this signed token is then passed to AWS to prove the identity of the runner.

\par In AWS, an OIDC provider is configured in IAM, trusting GitHub and allowing Role Assumptions for the worker as long as it proves it's from the Doclane repository. The IAM role assumed by the worker only has permissions to push images to the two ECR repositories. Nothing else.

\par Role Assumption in AWS is a concept that allows generating temporary credentials to perform actions in an AWS account, leveraging the STS (Secure Token Service) from AWS. Credentials are short-lived. Once the credentials expire they cannot be used anymore, and even if they somehow get leaked the damage would be minimal since the only thing they allow is pushing container images.

\section{Containerization}

\par Both the backend and frontend are packaged as Docker containers using multi-stage builds.

\subsection{Backend container}

\par The backend Dockerfile uses a two-stage build. The first stage uses the Go compiler to build a fully static binary with no external dependencies. The second stage uses a minimal distroless base image that contains no shell, no package manager, nothing except the binary and CA certificates. The final image is about 36 MB and runs as a non-root user.

\subsection{Frontend container}

\par The frontend Dockerfile also uses a multi-stage build. The first stage installs dependencies and runs the Next.js build. The second stage copies only the standalone output (a minimal server bundle) into a Node.js Alpine image. The final image is about 370 MB and runs as a non-root user.

\par The NEXT\_PUBLIC environment variables (Cognito pool ID, client ID, backend URL) are passed as build arguments since Next.js inlines them into the JavaScript bundle at build time. They are not secrets, they are public values that the browser needs to know.
